\chapter*{About this manual (English)}
\addcontentsline{toc}{chapter}{About this manual (English)}
\markboth{About this manual}{About this manual}

This is a complete, task-oriented handbook for \textbf{MobileCAD}, the Android front end of the FreeCAD engine.
It is written from the source code — the Kotlin/Compose interface, the Python bridge and the bundled FreeCAD
modules — and takes a reader who has never opened the app to a finished parametric part, an assembly, a
drawing, a finite-element run, toolpaths and a BIM model.

\begin{itemize}
  \item \textbf{Scope:} 17 chapters and 3 appendices, about 90 pages. Chapters 1--4 get a beginner productive
        (UI tour, gestures, first part); 5--7 cover the modelling core (sketch, PartDesign, Part); 8--14 cover
        the remaining workbenches (Draft, TechDraw, expressions, Assembly, FEM, CAM, BIM, mesh, points);
        15--17 cover files, efficiency and troubleshooting; the appendices hold a command reference, a bilingual
        glossary and an FAQ.
  \item \textbf{Accuracy:} commands, dialog rows and status messages are taken from the source (the command
        catalog is generated from FreeCAD's own sources), and limitations are stated where the port has them
        (e.g. no Netgen mesher, no CAM material-removal simulation, no DWG import).
  \item \textbf{Language:} the manual is written in Simplified Chinese because it accompanies the Chinese
        interface; every UI label is given as it appears in Chinese (the English command name is shown in
        parentheses on first use), and Appendix B maps both.
  \item \textbf{Screenshots:} from a Samsung Galaxy Tab S9+ (Android 16) running the Chinese build, plus the
        project's own marketing screenshots.
  \item \textbf{Building:} \file{cd manual \&\& xelatex main.tex} (run twice for the table of contents). It needs
        XeLaTeX with \file{ctex}, and a CJK font — the preamble uses macOS' Songti SC / Heiti SC; on other
        systems change the four \file{\textbackslash setCJK...font} lines in \file{main.tex}.
  \item \textbf{License:} GNU General Public License v3.0 or later, like the application. Contributions and
        corrections are welcome as issues or pull requests against \file{manual/}.
\end{itemize}
